% !TeX root = ../main.tex

\section{Related Work}
\label{sec:related}


\textbf{Token Reduction.} Prior work has shortened code inputs by removing or rewriting elements judged less important. DietCode, SlimCode, and LeanCode use model importance, program structure, or code naturalness to guide such simplification~\cite{dietcode2022,slimcode2024,leancode2025}; program-minimization studies similarly seek smaller inputs that preserve a prediction or expose model behavior~\cite{p2im2021,rabin2021,rabin2022}. These methods primarily target task-preserving or lossy reduction, rather than exact source-level recovery. Reversible methods instead define a compact representation that can be deterministically converted back to standard code. SimPy redesigns the Python surface grammar to remove layout and delimiter redundancy~\cite{simpy2024}, whereas TokenSugar mines recurring patterns and maps them to single-token shorthands~\cite{tokensugar2024}. Our work combines these two reduction levels and additionally studies whether a model-derived signal should decide where an otherwise valid shorthand is materialized.

This pattern-based direction is related to code-idiom mining, which extracts recurring AST, graph, or programmatic structures for abstraction, language design, or structured generation~\cite{allamanis2014,allamanis2018mining,sivaraman2022,iyer2019,shin2019}. NormCoRe differs in treating mined idioms as part of a reversible model-facing dialect: captures preserve occurrence-specific identifiers and expressions, while the decoder maintains the inverse mapping. The goal is therefore not only to find frequent idioms, but to select compact forms that remain valid, learnable, and recoverable.
This requirement is stricter than ordinary source simplification because a generated shorthand must remain well-formed under the same fixed registry and must be expanded without relying on a human repair step.

SimPy and TokenSugar also differ in how their compact forms are obtained. SimPy relies primarily on a human-designed grammar whose structural conventions are shared across inputs, whereas TokenSugar fixes a corpus-mined shorthand vocabulary and matching policy before model adaptation. This distinction creates a representation-learning question: a shorthand can be legal and token-efficient while still replacing cues that a model uses for prediction. NormCoRe therefore treats compression validity and occurrence selection as separate decisions. Normalization and Hybrid Idioms Shorthanding define the reversible candidate space; Conditional Attention decides whether an eligible occurrence should be materialized in shorthand form.

\textbf{Code Encoding and Model-Aware Representation.} Code models use subword tokenization and, in some cases, explicit identifier or program-structure signals~\cite{sennrich2016,karampatsis2020,wang2021,feng2020,guo2021,guo2022}. These methods improve code representation but do not primarily remove Python layout redundancy while providing a deterministic inverse codec. Interpretability studies investigate whether attention captures code syntax, data flow, or semantic structure~\cite{chen2022cat,wan2022,karmakar2021,bui2019,ma2022}. Model-aware prompt compression also shows that internal signals can guide input selection, although it targets general natural-language prompts and may discard information~\cite{llmlingua2023}. \tool extends this idea to reversible code representation: Conditional Attention uses contextual model evidence to select among legal shorthand occurrences without discarding the information required for restoration.

\textbf{Large Language Models for Generative Software Engineering.} Code completion, generation, testing, repository-level completion, and issue resolution require outputs that satisfy syntax, scope, interface, and execution constraints~\cite{raychev2014,chen2021codex,lemieux2023codamosa,liu2023repobench,jimenez2024swebench}. Repository-level tasks may additionally depend on definitions and usage examples distributed across files. A compact representation for these settings must therefore remain learnable and generatable while decoding deterministically to executable code; deleting a declaration or identifier can break cross-location dependencies even when a task's immediate prediction is preserved. Accordingly, \tool separates codec-side compression capacity from dialect learnability and evaluates whether model-informed occurrence selection improves the trade-off between reduction and downstream capability.
